
\documentclass{article}
\usepackage{colortbl}
\usepackage{makecell}
\usepackage{multirow}
\usepackage{supertabular}

\begin{document}

\newcounter{utterance}

\centering \large Interaction Transcript for game `referencegame', experiment `letter\_grids', episode 5 with reflexion{-}llama{-}3.1{-}8b{-}instruct.
\vspace{24pt}

{ \footnotesize  \setcounter{utterance}{1}
\setlength{\tabcolsep}{0pt}
\begin{supertabular}{c@{$\;$}|p{.15\linewidth}@{}p{.15\linewidth}p{.15\linewidth}p{.15\linewidth}p{.15\linewidth}p{.15\linewidth}}
    \# & $\;$A & \multicolumn{4}{c}{Game Master} & $\;\:$B\\
    \hline

    \theutterance \stepcounter{utterance}  
    & & \multicolumn{4}{p{0.6\linewidth}}{
        \cellcolor[rgb]{0.9,0.9,0.9}{
            \makecell[{{p{\linewidth}}}]{
                \texttt{\tiny{[P1$\langle$GM]}}
                \texttt{You have attempted this same task before and did not succeed. These are the plans you wrote after those attempts.} \\
\texttt{Trial \#1: New plan: I failed to distinguish the target grid from the distractor grids because my description was too general and didn't provide enough unique features. The target grid has a unique pattern of X's and □'s, but my previous description only mentioned the positions of the □'s in the second column, which is also present in the distractor grids. Next time, I will focus on describing the unique arrangement of X's and □'s in the target grid, such as the fact that the top and bottom rows have X's in every column, while the middle rows have X's only in the first and last columns. I will also try to find a more specific and detailed description of the □'s in the target grid, such as their positions in the second, third, and fourth columns.} \\
\\ 
\texttt{Trial \#2: New plan: In my previous attempt, I focused on describing the unique arrangement of X's and □'s in the target grid, but I failed to distinguish it from the distractor grids. Upon re{-}examining the grids, I noticed that the target grid has a unique feature: the □'s in the second, third, and fourth columns of the middle rows are all adjacent to each other, forming a continuous block of empty cells. This feature is not present in the distractor grids. Next time, I will focus on describing this unique feature, such as "The grid that has a continuous block of □'s in the second, third, and fourth columns of the middle rows."} \\
\\ 
\\ 
\texttt{You are given three grids, where each of them is 5 by 5 in size.} \\
\texttt{Grids have empty cells marked with "▢" and filled cells marked with "X".} \\
\texttt{Your task is to generate a referring expression that best describes the target grid while distinguishing it from the two other distractor grids.} \\
\texttt{The first grid is the target grid, and the following two grids are the distractors.} \\
\\ 
\texttt{Target grid:} \\
\\ 
\texttt{X X X X X} \\
\texttt{X □ □ □ □} \\
\texttt{X □ □ □ □} \\
\texttt{X □ □ □ □} \\
\texttt{X X X X X} \\
\\ 
\texttt{Distractor grid 1:} \\
\\ 
\texttt{X X X X X} \\
\texttt{X □ □ □ □} \\
\texttt{X X X X X} \\
\texttt{X □ □ □ □} \\
\texttt{X X X X X} \\
\\ 
\texttt{Distractor grid 2:} \\
\\ 
\texttt{X X X X X} \\
\texttt{X □ X □ X} \\
\texttt{X X X X X} \\
\texttt{X □ X □ X} \\
\texttt{X X X X X} \\
\\ 
\texttt{Instruction: Describe the target grid.} \\
\texttt{Generate the referring expression starting with the tag "Expression: " for the given target grid. Omit any other text.} \\
            }
        }
    }
    & \\ \\

    \theutterance \stepcounter{utterance}  
    & \multicolumn{4}{p{0.6\linewidth}}{
        \cellcolor[rgb]{0.9,0.9,0.9}{
            \makecell[{{p{\linewidth}}}]{
                \texttt{\tiny{[P1$\rangle$GM]}}
                \texttt{Expression: The grid that has a continuous block of □'s in the second, third, and fourth columns of the middle rows.} \\
            }
        }
    }
    & & \\ \\

    \theutterance \stepcounter{utterance}  
    & & & \multicolumn{2}{p{0.3\linewidth}}{
        \cellcolor[rgb]{0.9,0.9,0.9}{
            \makecell[{{p{\linewidth}}}]{
                \texttt{\tiny{[GM$|$GM]}}
                \texttt{Expression: The grid that has a continuous block of □'s in the second, third, and fourth columns of the middle rows.} \\
            }
        }
    }
    & & \\ \\

    \theutterance \stepcounter{utterance}  
    & & \multicolumn{4}{p{0.6\linewidth}}{
        \cellcolor[rgb]{0.9,0.9,0.9}{
            \makecell[{{p{\linewidth}}}]{
                \texttt{\tiny{[GM$\rangle$P2]}}
                \texttt{You have attempted this same task before and did not succeed. These are the plans you wrote after those attempts.} \\
\texttt{Trial \#1: New plan: I selected the second grid because it seemed to match the description of having four rows of X's at the top and bottom, and three rows of □'s in the middle. However, upon re{-}reading the description, I realize that I misinterpreted the phrase "with the □'s in the second column." The correct interpretation is that the □'s should be in the middle rows, not necessarily in the second column. Next time, I will pay closer attention to the description and look for a grid that matches the correct interpretation. I will also consider the possibility that the description might be referring to the arrangement of □'s in a different column.} \\
\\ 
\texttt{Trial \#2: New plan: I selected the second grid because it seemed to match the description of having four rows of X's at the top and bottom, and three rows of □'s in the middle. However, upon re{-}reading the description, I realize that I misinterpreted the phrase "with the □'s in the second, third, and fourth columns of the middle rows." The correct interpretation is that the □'s should be in the middle rows, and specifically in the second, third, and fourth columns. Next time, I will pay closer attention to the description and look for a grid that matches the correct interpretation. I will also consider the possibility that the description might be referring to the arrangement of □'s in a different column, and I will look for a grid that has □'s in the second, third, and fourth columns of the middle rows, regardless of the column in which they appear in the top and bottom rows.} \\
\\ 
\\ 
\texttt{You are given three grids, where each of them is 5 by 5 in size.} \\
\texttt{Grids have empty cells marked with "▢" and filled cells marked with "X".} \\
\texttt{You are also given a referring expression that describes one of the given grids.} \\
\texttt{Your task is to select the grid that matches the given referring expression.} \\
\texttt{Generate only the number (in text) of the grid that matches the given expression by selecting first, second, or third.} \\
\\ 
\texttt{First:} \\
\\ 
\texttt{X X X X X} \\
\texttt{X □ X □ X} \\
\texttt{X X X X X} \\
\texttt{X □ X □ X} \\
\texttt{X X X X X} \\
\\ 
\texttt{Second:} \\
\\ 
\texttt{X X X X X} \\
\texttt{X □ □ □ □} \\
\texttt{X X X X X} \\
\texttt{X □ □ □ □} \\
\texttt{X X X X X} \\
\\ 
\texttt{Third:} \\
\\ 
\texttt{X X X X X} \\
\texttt{X □ □ □ □} \\
\texttt{X □ □ □ □} \\
\texttt{X □ □ □ □} \\
\texttt{X X X X X} \\
\\ 
\texttt{Expression: The grid that has a continuous block of □'s in the second, third, and fourth columns of the middle rows.} \\
\texttt{Question: Which grid does the expression refer to?} \\
\texttt{Start with the tag "Answer: ", followed by your selection. Omit any other text.} \\
            }
        }
    }
    & \\ \\

    \theutterance \stepcounter{utterance}  
    & & & \multicolumn{4}{p{0.6\linewidth}}{
        \cellcolor[rgb]{0.9,0.9,0.9}{
            \makecell[{{p{\linewidth}}}]{
                \texttt{\tiny{[GM$\langle$P2]}}
                \texttt{Answer: third} \\
            }
        }
    }
     \\ \\

    \theutterance \stepcounter{utterance}  
    & & & \multicolumn{2}{p{0.3\linewidth}}{
        \cellcolor[rgb]{0.9,0.9,0.9}{
            \makecell[{{p{\linewidth}}}]{
                \texttt{\tiny{[GM$|$GM]}}
                \texttt{third} \\
            }
        }
    }
    & & \\ \\

\end{supertabular}
}

\end{document}
